\documentclass[12pt,a4paper]{article}

% ---- Page Layout (Spec: 3.5cm left, 2.5cm top, 1.25cm right, 2.0cm bottom) ----
\usepackage[a4paper, left=3.5cm, right=1.25cm, top=2.5cm, bottom=2.0cm]{geometry}

% ---- Fonts & Spacing (Spec: Times New Roman, 12pt, 1.5 spacing, justified) ----
\usepackage{setspace}
\usepackage{newtxtext, newtxmath}
\usepackage[T1]{fontenc}
\usepackage[utf8]{inputenc}

% ---- Math, Tables, Figures ----
\usepackage{amsmath}
\usepackage{amssymb}
\usepackage{graphicx}
\usepackage{booktabs}
\usepackage{array}
\usepackage{caption}

% ---- Hyperlinks & TOC ----
\usepackage[hidelinks]{hyperref}
\usepackage{float}

% ---- Caption placement: Figure captions BOTTOM (default), Table captions TOP ----
\captionsetup[table]{position=above}
\captionsetup[figure]{position=below}

% ---- Paragraph formatting (justified, no extra indent after sections) ----
\usepackage{parskip}

% ---- Page style ----
\pagestyle{plain}

\onehalfspacing

\begin{document}

%================================================================
%  TITLE PAGE  (BUG-04 fix)
%================================================================
\begin{titlepage}
\centering
\vspace*{2cm}
{\LARGE\bfseries Adaptive Semantic-Velocity Chunking\\[6pt]
with Dynamic-Relevant RAG (DR-RAG)\par}
\vspace{1.5cm}
{\large A Major Project Report\\[4pt]
Submitted in Partial Fulfillment of the Requirements\\[4pt]
for the Degree of Bachelor of Technology\par}
\vspace{2cm}
{\Large\itshape Rajeev Gupta\par}
\vfill
{\large Department of Computer Science and Engineering\\[6pt]
\today\par}
\end{titlepage}

%================================================================
%  INDEX / TABLE OF CONTENTS  (Spec §3: Index after title page)
%================================================================
\tableofcontents
\newpage

%================================================================
%  ABSTRACT  (Spec: should not exceed 1 page)
%================================================================
\begin{abstract}

Retrieval-Augmented Generation (RAG) systems face fundamental limitations when processing
long, structured documents such as legal contracts and academic papers. Standard fixed-size
chunking disrupts semantic coherence, while single-stage retrieval fails to support
multi-hop reasoning across weakly linked evidence.

This work proposes \textbf{Adaptive Semantic-Velocity Chunking with Dynamic-Relevant RAG
(DR-RAG)}, a two-phase framework that addresses both challenges. Phase~I introduces
\textbf{Context-Augmented Agentic Chunking (CAAC)}, which replaces static chunking with
Small Language Model (SLM)-guided semantic boundaries. Each chunk is enriched with a
rolling summary of preceding context, ensuring preservation of both local meaning and
global document flow. Semantic velocity, defined as embedding drift between consecutive
chunks, dynamically adjusts chunk size and summary length.

Phase~II introduces a \textbf{classifier-routed two-stage retrieval mechanism}. The first
stage retrieves anchor documents using cosine similarity. For complex queries, a second
stage generates expanded queries from intermediate facts to retrieve additional supporting
evidence. The final response is generated using the union of anchor and evidence sets.

The proposed framework is designed to improve recall, contextual coherence, multi-hop
reasoning accuracy, and hallucination resistance compared to conventional RAG systems.
Experimental evaluation is planned using HotpotQA~\cite{yang2018hotpotqa} and
CUAD~\cite{hendrycks2021cuad} benchmarks.

\end{abstract}

\newpage

%================================================================
%  MOTIVATION  (Spec: should not exceed 1 page)
%================================================================
\section{Motivation}

Retrieval-Augmented Generation (RAG) has emerged as the dominant paradigm for
grounding large language model (LLM) responses in external, verifiable knowledge.
By retrieving relevant document passages at inference time, RAG systems reduce
hallucination and extend the effective knowledge horizon of LLMs beyond their
training cut-off. Despite this promise, current RAG pipelines exhibit critical
structural weaknesses when applied to long, domain-specific documents such as
legal contracts, academic research papers, and technical standards. Three
fundamental limitations motivate the design of the proposed DR-RAG framework.

\begin{description}
\item[Limitation 1: Semantic Fragmentation due to Fixed-Size Chunking.]
The overwhelming majority of deployed RAG systems partition source documents
into fixed-size token windows of 256--512 tokens with optional overlap.
This strategy is computationally convenient but semantically blind. In legal
documents, a single obligation clause may span several hundred words across
multiple sentences; splitting it at a token boundary severs the conditional
relationship between its premise and its consequence. Similarly, in academic
papers, a mathematical proposition and its proof are logically inseparable,
yet fixed-size chunking routinely places them in different chunks. The
retriever then returns a fragment that contains only part of the relevant
reasoning, causing the generator to produce incomplete or factually incorrect
answers. Overlap-based mitigation inflates the vector index size without
resolving the underlying boundary problem.

\item[Limitation 2: Context Decay in Long Documents.]
Even when individual chunks are semantically coherent, they are stored and
retrieved in isolation. Definitions, assumptions, and conditions established
in the introduction or background sections of a document are invisible to a
chunk retrieved from a later section. This \emph{context decay} is especially
damaging in legal and scientific texts, where a clause in section~8 may be
legally conditioned on a definition given in section~2, or where an
experimental result in section~4 is only interpretable in light of the
methodology described in section~3. No existing chunking framework provides
a mechanism to carry forward this accumulated contextual state into each
retrieved unit.

\item[Limitation 3: Failure of Single-Stage Retrieval for Multi-Hop Reasoning.]
Many practical queries over structured documents require multi-hop reasoning:
the answer cannot be found in any single passage but must be assembled from
evidence distributed across two or more weakly linked sections. Single-stage
dense retrieval, which ranks passages by cosine similarity to the original
query, retrieves only the most superficially similar chunk and misses the
bridging evidence required to complete the reasoning chain. While iterative
and graph-based approaches such as GraphRAG and multi-agent retrieval systems
partially address this, they do so at the cost of constructing and maintaining
expensive knowledge graphs or issuing multiple sequential LLM calls, making
them impractical for latency-sensitive production environments.
\end{description}

\newpage
%================================================================
%  LITERATURE REVIEW
%================================================================
\section{Literature Review}

The design of an effective RAG system depends critically on two components:
how documents are segmented into retrievable units (chunking), and how
relevant units are identified at query time (retrieval). This review surveys
the key chunking strategies proposed in the literature, examines existing
retrieval architectures, and identifies the gaps that motivate the proposed
DR-RAG framework.

%----------------------------------------------------------------
\subsection{Chunking Strategies in RAG Systems}
%----------------------------------------------------------------

Chunking is the process of partitioning a source document into a collection
of retrievable units. The granularity, coherence, and contextual completeness
of these units directly determine the upper bound on retrieval quality.
Table~\ref{tab:lit_review} summarises the principal chunking methods, their
mechanisms, and their trade-offs.

% Spec §6: Table caption at TOP
\begin{table}[!htbp]
\centering
\caption{Comparison of Existing Chunking Methods for RAG Systems}
\label{tab:lit_review}
\renewcommand{\arraystretch}{1.3}
\begin{tabular*}{\textwidth}{@{\extracolsep{\fill}} p{2.8cm} p{4.5cm} p{3.0cm} p{3.0cm}}
\toprule
\textbf{Method} & \textbf{Description} & \textbf{Pros} & \textbf{Cons} \\
\midrule
Fixed-Size
  & Splits document into windows of $n$ tokens with optional overlap
  & Simple; fast; predictable index size
  & Breaks semantic structure; mid-sentence splits common \\[2pt]

Sentence-Based
  & Uses punctuation or NLP sentence detectors to set boundaries
  & Improves local coherence over fixed-size
  & Ignores paragraph and section hierarchy \\[2pt]

Recursive
  & Applies a hierarchy of delimiters (paragraph, sentence, word)
  & Flexible granularity; structure-aware
  & Boundaries remain arbitrary; no semantic validation \\[2pt]

Semantic
  & Computes embedding similarity between adjacent segments; splits on high drift
  & Produces semantically coherent chunks
  & Computationally expensive; sensitive to embedding quality \\[2pt]

Proposition-Based
  & Decomposes text into atomic, self-contained factual propositions
  & High retrieval precision; minimal noise per chunk
  & Very slow; requires strong generative model; costly at scale \\[2pt]

Page-Level
  & Treats each document page as a single chunk
  & Preserves physical layout and visual structure
  & Chunks too large; poor retrieval precision; cross-page context lost \\[2pt]

Hierarchical (RAPTOR)
  & Builds a tree of progressively summarised clusters via recursive abstraction
  & Supports multi-granularity retrieval
  & Complex indexing pipeline; high memory overhead; slow to build \\

\bottomrule
\end{tabular*}
\end{table}

\subsubsection{Fixed-Size and Sentence-Based Chunking}

Fixed-size chunking~\cite{lewis2020rag} is the most widely adopted approach
due to its simplicity and deterministic behaviour. A sliding window of a
fixed number of tokens (typically 256--512) is moved across the document,
and each window becomes one chunk. While fast and easy to implement, this
method is agnostic to sentence, clause, and paragraph boundaries. Legal
clauses or theorem proofs that span more than one window are split mid-way,
causing retrieval to return semantically incomplete chunks. Overlapping
windows partially mitigate this but inflate index size and introduce
redundancy.

Sentence-based chunking improves local coherence by using NLP tools (e.g.,
spaCy, NLTK punkt tokenizer) to detect sentence boundaries. However, it
treats each sentence in isolation and ignores the structural hierarchy of the
document. A legal sub-clause spread across three sentences may be split
across two chunks, losing the conditional relationship between them.

\subsubsection{Recursive and Semantic Chunking}

Recursive chunking applies a priority-ordered list of delimiters
(e.g., \texttt{\textbackslash n\textbackslash n}, \texttt{\textbackslash n},
full stops) and recursively splits segments that exceed a maximum token
limit. While more structure-aware than fixed-size methods, the boundaries
produced are still syntactic rather than semantic, and the recursion depth
is a manually tuned hyperparameter.

Semantic chunking~\cite{gao2024survey} represents a significant improvement.
It embeds each sentence and computes the cosine similarity between adjacent
segments. A chunk boundary is inserted wherever the similarity drops below
a predefined threshold, indicating a topic shift. This approach produces
topically coherent chunks but requires an embedding model pass over the
entire document at ingestion time, making it expensive for large corpora.
It also lacks any mechanism to carry forward contextual information from
earlier chunks.

\subsubsection{Proposition-Based Chunking}

Proposition-based chunking~\cite{chen2023dense} decomposes text into atomic,
self-contained factual statements. Each proposition is a minimal claim that
can be assessed independently. This maximises retrieval precision since each
chunk is informationally pure. However, the extraction requires a capable
instruction-following language model, making this method approximately
10--20$\times$ slower than fixed-size chunking. It is impractical for
real-time or large-scale indexing pipelines.

\subsubsection{Page-Level and Hierarchical (RAPTOR) Chunking}

Page-level chunking is used in document-image pipelines where the physical
layout carries meaning (e.g., scanned PDFs). Each page is treated as a
single chunk. While this preserves visual structure, pages are too coarse
as retrieval units and frequently span multiple unrelated topics.

RAPTOR~\cite{sarthi2024raptor} constructs a tree of summaries by recursively
clustering and summarising document segments bottom-up. At query time,
retrieval can occur at multiple levels of the tree, enabling both fine-grained
and abstract matching. Despite its theoretical elegance, RAPTOR requires
building and maintaining a complex tree index, incurs high memory overhead,
and its summarisation quality degrades for highly technical documents such
as legal contracts.

%----------------------------------------------------------------
\subsection{Retrieval Architectures}
%----------------------------------------------------------------

\subsubsection{Single-Stage Dense Retrieval}

Dense Passage Retrieval (DPR)~\cite{karpukhin2020dpr} introduced bi-encoder
models that independently encode queries and passages into a shared embedding
space, enabling efficient maximum inner-product search (MIPS). DPR achieves
strong performance on factoid open-domain QA benchmarks but is limited to
single-hop questions where the answer is contained within a single passage.
For multi-hop questions requiring evidence from multiple documents, its
recall degrades substantially.

\subsubsection{Multi-Stage and Iterative Retrieval}

To address multi-hop reasoning, several iterative retrieval approaches have
been proposed. REPLUG~\cite{yu2023replug} treats the retrieval corpus as a
plug-in module for black-box LLMs and uses ensemble scoring over multiple
retrieved documents. Atlas~\cite{izacard2022atlas} jointly fine-tunes a
retriever and a reader using few-shot supervision, achieving strong results
on knowledge-intensive tasks. However, both methods require repeated forward
passes through large models and scale poorly with corpus size.

\subsubsection{Graph-Based Retrieval}

GraphRAG~\cite{edge2024graphrag} constructs a knowledge graph from the
document corpus and traverses entity-relation edges to answer complex
queries. While it excels at reasoning over interconnected entities, graph
construction is expensive and requires structured entity extraction pipelines.
The latency of graph traversal also makes it unsuitable for real-time
applications.

%----------------------------------------------------------------
\subsection{Research Gaps}
%----------------------------------------------------------------

A systematic analysis of the above literature reveals four critical gaps
that existing methods fail to address simultaneously:

\begin{description}
\item[1. Absence of Document-Type-Aware Chunking.]
All surveyed chunking methods treat documents uniformly, applying the same
splitting strategy regardless of whether the input is a legal contract, a
scientific paper, or a news article. Legal documents are structured around
hierarchical clause trees, while academic papers follow IMRaD conventions.
A chunking strategy that is agnostic to these structural patterns will
systematically fragment semantically complete units.

\item[2. Context Decay in Long Documents.]
Fixed-size, sentence-based, and even semantic chunking methods discard all
preceding content when constructing each chunk. In long documents, this
leads to \emph{context decay}: later chunks lose access to definitions,
assumptions, or conditions established in earlier sections. No existing
chunking framework incorporates a rolling contextual memory mechanism to
counteract this decay.

\item[3. Expensive Multi-Hop Retrieval.]
Methods capable of multi-hop reasoning (RAPTOR, GraphRAG, iterative
retrieval) impose high computational and latency costs. GraphRAG requires
full graph construction before any query can be answered. Iterative
retrievers require multiple sequential LLM calls. These costs are
prohibitive for production deployments requiring sub-second response times.

\item[4. Absence of Formal Chunking Quality Metrics.]
The chunking literature lacks a formal, quantitative metric for evaluating
whether a proposed boundary is semantically appropriate. Existing methods
rely on downstream task performance (e.g., QA accuracy) as a proxy, making
it difficult to diagnose chunking failures independently of retrieval and
generation. A dedicated metric such as \emph{semantic velocity} --- measuring
embedding drift at chunk boundaries --- would enable principled, task-agnostic
evaluation of chunking quality.
\end{description}

The proposed DR-RAG framework directly targets all four gaps through
(i)~SLM-guided, structure-aware boundary detection; (ii)~rolling summary
augmentation for context preservation; (iii)~a lightweight classifier-routed
two-stage retrieval pipeline; and (iv)~semantic velocity as a formal chunking
quality metric. The relationship between each gap and its corresponding
system component is summarised in Table~\ref{tab:gap_mapping}.

% Spec §6: Table caption at TOP
\begin{table}[!htbp]
\centering
\caption{Mapping of Research Gaps to Proposed System Components}
\label{tab:gap_mapping}
\renewcommand{\arraystretch}{1.4}
\begin{tabular*}{\textwidth}{@{\extracolsep{\fill}} p{0.6cm} p{4.8cm} p{7.0cm}}
\toprule
\textbf{\#} & \textbf{Research Gap} & \textbf{DR-RAG Component} \\
\midrule
G1 & No document-type-aware chunking
   & SLM-guided boundary detector with structure prompts (CAAC \S1.2) \\
G2 & Context decay in long documents
   & Rolling summary augmentation with velocity-regulated budget (CAAC \S1.4) \\
G3 & Expensive multi-hop retrieval
   & Classifier-routed two-stage retrieval; no graph required (DR-RAG \S2.2--2.4) \\
G4 & No formal chunking quality metric
   & Semantic velocity: embedding drift between adjacent chunks (CAAC \S1.3) \\
\bottomrule
\end{tabular*}
\end{table}

\newpage
%================================================================
%  METHODOLOGY
%================================================================
\section{Methodology}

\subsection{Objective}

The primary objective of this work is to design an efficient and scalable
Retrieval-Augmented Generation (RAG) framework that addresses two fundamental
limitations of existing systems:
    (i)~preservation of semantic document structure through adaptive, context-aware
chunking, and
    (ii)~accurate multi-hop reasoning via a lightweight classifier-routed two-stage
retrieval mechanism.

The overall system is referred to as \textbf{DR-RAG} (Dynamic-Relevant RAG) and
consists of two sequential phases: Phase~I handles document ingestion and
chunk construction, while Phase~II handles query processing and response
generation. The end-to-end pipeline takes an input document corpus and a
user query and produces a generated response.

% \begin{figure}[!htbp]
%     \centering
%     \IfFileExists{high level architecture digram.png}{%
%         \includegraphics[width=\linewidth]{high level architecture digram.png}%
%     }{%
%         \fbox{\parbox{0.9\linewidth}{\centering\vspace{2cm}
%         \textit{[Architecture diagram placeholder --- insert high level architecture digram.png here]}
%         \vspace{2cm}}}%
%     }
%     \caption{System Architecture of DR-RAG}
%     \label{fig:architecture}
% \end{figure}


%-------- Phase I --------%
\subsection{Phase I: Context-Augmented Agentic Chunking (CAAC)}

\subsubsection{Overview}

Standard chunking methods partition a document into fixed-size token windows
without regard for semantic or structural boundaries, causing context
fragmentation in long structured documents such as legal contracts and research
papers. Phase~I addresses this by introducing \textbf{Context-Augmented Agentic
Chunking (CAAC)}: a semantics-driven strategy guided by a Small Language
Model~(SLM) and regulated by a novel metric called \emph{semantic velocity}.

\begin{figure}[!htbp]
    \centering
    \IfFileExists{Phase 1 architecture diagram.png}{%
        \includegraphics[width=\linewidth]{Phase 1 architecture diagram.png}%
    }{%
        \fbox{\parbox{0.9\linewidth}{\centering\vspace{2cm}
        \textit{[Architecture diagram placeholder --- insert Phase 1 architecture diagram.png here]}
        \vspace{2cm}}}%
    }
    \caption{Phase I: Context-Augmented Agentic Chunking (CAAC) Architecture}
    \label{fig:phase1_architecture}
\end{figure}

\subsubsection{Two-Pass Document Segmentation}

To resolve the circular dependency between chunk boundaries and semantic
velocity (which requires chunk embeddings that in turn depend on chunk
boundaries), DR-RAG employs a \textbf{two-pass hierarchical segmentation
process}.

\paragraph{Pass~1: Preliminary Sentence Embedding.}
The document is first segmented into individual sentences. Each sentence is
embedded independently using the embedding model, and the distance between
adjacent sentence embeddings is computed. These sentence-level drift values
are used to generate preliminary chunk boundaries via the SLM boundary
detector.

\paragraph{Pass~2: Adaptive Chunk Refinement.}
The preliminary chunks are then merged or split dynamically based on their
semantic drift values:
\begin{itemize}
  \item If drift is \textbf{low} (below a lower threshold), adjacent chunks
    are \textbf{merged} since they cover the same topic.
  \item If drift is \textbf{moderate} (between thresholds), the chunk
    boundary is \textbf{retained} as-is.
  \item If drift is \textbf{high} (above an upper threshold), the chunk is
    further \textbf{split} to isolate the topic change.
\end{itemize}
This two-pass design eliminates the circular dependency because sentence-level
embeddings are computed \emph{before} adaptive chunk boundaries are formed.
Each final chunk corresponds to a logically complete unit of information
(e.g., a legal clause, a theorem, or a methodology paragraph).

\subsubsection{Boundary Detector Training Data}

The SLM boundary detector is fine-tuned using a semi-supervised dataset
constructed from:
\begin{itemize}
  \item legal contracts (CUAD~\cite{hendrycks2021cuad}),
  \item academic papers (arXiv / PubMed),
  \item manually annotated semantic transitions.
\end{itemize}
Weak supervision labels~\cite{ratner2017snorkel} are generated using heading
transitions, paragraph topic shifts, and citation density changes. Human
verification is applied to a sampled subset to ensure label quality.

\subsubsection{Semantic Velocity}

\emph{Semantic velocity} quantifies the magnitude of meaning shift between
two consecutive chunks. It is computed as the Euclidean distance between the
embedding vectors of adjacent chunks --- a higher value indicates a larger
topic shift. Semantic velocity drives two dynamic control functions:

\begin{enumerate}
  \item \textbf{Adaptive chunk sizing:} Low semantic velocity causes adjacent
    segments to be merged into a larger coherent chunk; high velocity enforces
    an immediate hard boundary.
  \item \textbf{Summary length regulation:} Higher semantic velocity triggers
    a longer rolling summary to bridge the contextual gap introduced by the
    topic shift.
\end{enumerate}

\subsubsection{Context-Augmented Chunk Construction}

Each raw chunk is augmented with a rolling summary generated from the
\textbf{previous three chunks}. A window size of three was selected because
a single preceding chunk provides insufficient long-range context, while
more than five introduces excessive redundancy and token overhead; empirical
pilot observations confirmed that three provides a practical balance between
coherence and efficiency. The rolling summary is prepended (concatenated) to
the original chunk content, producing the final \emph{augmented chunk}.

\subsubsection{Token Budget Scaling}

The rolling-summary token budget scales linearly with semantic velocity.
The budget ranges from a minimum of 32~tokens (for low drift) to a maximum
of 128~tokens (for high drift). This means smooth topic continuity receives
short summaries, while large topic shifts receive longer summaries to bridge
the contextual gap.

\subsubsection{Context Retention Score}

To evaluate whether a chunk remains semantically understandable in isolation,
we define the \textbf{Context Retention Score (CRS)}. CRS is calculated as
the ratio of the cosine similarity between the chunk and its original
surrounding document context~\cite{reimers2019sentencebert} to the chunk's
token length. A higher CRS indicates that the chunk preserves more of its
original meaning without depending on external context.

\subsubsection{Embedding and Indexing}

Each augmented chunk is encoded into a dense vector using the
BGE-M3~\cite{chen2024bge} bi-encoder embedding model. All vectors are
stored in a vector database (FAISS~\cite{johnson2019faiss} /
Qdrant~\cite{qdrant2023}) and indexed with Hierarchical Navigable Small
World~(HNSW) graphs~\cite{malkov2020hnsw} for efficient approximate
nearest-neighbour search at query time.


%-------- Phase II --------%
\subsection{Phase II: Dynamic-Relevant RAG (DR-RAG)}

\subsubsection{Overview}

Phase~II defines the query-time pipeline. A lightweight classifier first
determines the structural complexity of the incoming query. Simple factoid
queries are answered with single-stage anchor retrieval. Complex multi-hop
queries additionally trigger an evidence-mining stage, where intermediate
facts extracted from anchor chunks are used to formulate a secondary query.
The final answer is generated conditioned on the union of both retrieval sets.

\begin{figure}[!htbp]
    \centering
    \IfFileExists{phase 2 architecture diagram.png}{%
        \includegraphics[width=\linewidth]{phase 2 architecture diagram.png}%
    }{%
        \fbox{\parbox{0.9\linewidth}{\centering\vspace{2cm}
        \textit{[Architecture diagram placeholder --- insert phase 2 architecture diagram.png here]}
        \vspace{2cm}}}%
    }
    \caption{Phase II: Dynamic-Relevant RAG (DR-RAG) Architecture}
    \label{fig:phase2_architecture}
\end{figure}

\subsubsection{Step 1 --- Query Encoding and Routing}

The incoming natural language query is encoded into a dense vector and
passed to a \textbf{query routing classifier}. The classifier is implemented
as a lightweight transformer-based binary classifier fine-tuned on labeled
query-complexity data. It uses the following features:
\begin{itemize}
  \item the query's embedding vector,
  \item dependency-tree depth (syntactic complexity),
  \item conjunction frequency,
  \item semantic entropy score.
\end{itemize}
Based on these features, the classifier routes each query as either
\textbf{single-hop} (answerable from one passage) or \textbf{multi-hop}
(requires evidence from multiple passages).

\paragraph{Training Data.}
The classifier is trained using HotpotQA~\cite{yang2018hotpotqa} multi-hop
labels, Natural Questions~\cite{kwiatkowski2019nq} single-hop queries, and
manually labeled legal QA pairs from CUAD~\cite{hendrycks2021cuad}.

\paragraph{Confidence Fallback.}
If the classifier confidence falls in the ambiguous range (between 45\% and
55\%), the system defaults to two-stage retrieval to maximise recall.

\subsubsection{Step 2 --- Anchor Retrieval}

For all query types, the top-\textit{k} most semantically similar chunks are
retrieved from the vector index using cosine similarity, forming the
\emph{anchor set}. For single-hop queries, the anchor set serves directly as
the final context. For multi-hop queries, it provides the seed for evidence
mining.

\subsubsection{Step 3 --- Evidence Mining (Multi-Hop Only)}

\paragraph{(a) Intermediate Fact Extraction.}
The fact extraction step uses \textbf{SLM-guided structured extraction}.
Each anchor chunk is transformed into atomic propositions using prompt-based
information extraction. The extraction process combines dependency parsing,
Open Information Extraction~(OpenIE)~\cite{stanovsky2018openie}, and SLM
prompt refinement. For example, given the chunk \textit{``Clause~8 applies
only when prior written consent is obtained,''} the extracted proposition
would be the triple: (Clause~8, requires, WrittenConsent). This hybrid
design reduces hallucinated facts while preserving semantic relationships.

\paragraph{(b) Mining Query Synthesis.}
The extracted facts are combined with the original query to synthesise an
expanded secondary query that targets the missing evidence.

\paragraph{(c) Evidence Retrieval.}
The expanded query is encoded and used to retrieve additional evidence chunks
from the vector index, forming the \emph{evidence set}. Retrieved chunks
that overlap with the anchor set are excluded to ensure complementary
coverage.

\subsubsection{Context Assembly and Response Generation}

Anchor and evidence chunks are deduplicated using cosine similarity with a
threshold of 0.95: if two chunks have a similarity score at or above this
value, they are considered duplicates and the one with the higher retrieval
relevance score is retained.

The language model then generates the final response autoregressively,
conditioned on the query and the assembled context. The prompt passed to
the generation model follows the structure: \textit{[System Instruction]
+ [Assembled Context] + [Query]}.

%-------- Ablation Study --------%
\subsection{Ablation Study}

To isolate the contribution of each DR-RAG component, four ablation
configurations will be evaluated:

\begin{table}[!htbp]
\centering
\caption{Ablation Study Configurations}
\label{tab:ablation}
\renewcommand{\arraystretch}{1.3}
\begin{tabular*}{\textwidth}{@{\extracolsep{\fill}} c p{10cm}}
\toprule
\textbf{Config} & \textbf{Removed Component} \\
\midrule
A1 & No rolling summaries \\
A2 & Fixed-size chunking instead of CAAC \\
A3 & No query routing (all queries use two-stage) \\
A4 & No evidence mining stage \\
\bottomrule
\end{tabular*}
\end{table}

Metrics compared across all configurations:
\begin{itemize}
  \item Recall@$K$
  \item Exact Match (EM)
  \item Hallucination Rate
  \item Latency
  \item Context Retention Score (CRS)
\end{itemize}

This ablation determines which architectural component contributes most
significantly to performance.

\newpage

%================================================================
%  BIBLIOGRAPHY / REFERENCES  (Spec §7: IEEE format)
%================================================================
\bibliographystyle{IEEEtran}
\bibliography{references}

\end{document}